\documentclass[11pt,a4paper,twoside]{article}

\usepackage[T1]{fontenc}
\usepackage{lmodern}
\usepackage{microtype}
\usepackage[a4paper,margin=29mm,headheight=14pt]{geometry}
\usepackage{amsmath,amssymb,mathtools,amsthm}
\usepackage{booktabs,tabularx,longtable,array}
\usepackage{enumitem}
\usepackage[dvipsnames]{xcolor}
\usepackage{fancyhdr}
\usepackage{listings}
\usepackage{url}
\usepackage{hyperref}
\usepackage[nameinlink,capitalize,noabbrev]{cleveref}

% Generated by paper/finalize_publication.ps1 from the accepted bundle-v2 audit.
% Do not edit this file by hand.
\newif\ifcertificateaccepted
\certificateacceptedtrue

\newcommand{\CertifiedConstant}{\(\frac{9}{20}=0.45\)}
\newcommand{\CertificateBundleFormat}{berry-esseen-045-certificate-bundle-v2}
\newcommand{\CertificateScopeStart}{2}
\newcommand{\CertificateScopeEnd}{8123}
\newcommand{\CertificateTarget}{9/20}
\newcommand{\CertificatePythonFlint}{0.9.0}
\newcommand{\CertificateFlint}{3.6.0}
\newcommand{\CertificateDate}{2026-08-06T20:09:07Z}
\newcommand{\CertificateTotalCells}{886998}
\newcommand{\CertificateWorstUpper}{0.449998514435452506681043349522}
\newcommand{\CertificateWorstUpperExact}{5649372911768414488969095178501869845574605158618979945689/12554203470773361527671578846415332832204710888928069025792}
\newcommand{\CertificateWorstLocation}{\(n=13\), cell \(21\), on \(\rho\in[304162949/301568000,152143259/150784000]\)}
\newcommand{\CertificateElapsedSeconds}{38558.206}
\newcommand{\CertificateSetDigest}{df1986ed406dc14d1ab071c4812adeadc1ad6671dff31cdbdf49f06415b70297}
\newcommand{\CertificateManifestDigest}{0fb539f0341492181a2eb23d0b2a512c\allowbreak{}ab50a819a2582c955f24cb5dd22a24f5}
\newcommand{\CertificateCheckerDigest}{22594f73588800cbdb785b563d570b52f180688602e2761ba1ff3ef2c729a4a6}
\newcommand{\CertificateCoreDigest}{6aa2f8228528e886c03e0192a44690f61b6bec1335fb6d083fe0a2a2644702b9}
\newcommand{\CertificateRuntime}{\code{Linux-6.6.87.2-microsoft-standard-WSL2-x86\_64-with-glibc2.36}}


\ifcertificateaccepted
  \newcommand{\PaperTitleText}{A 0.45 Berry--Esseen Bound}
  \newcommand{\PaperRunningTitle}{A 0.45 Berry--Esseen bound}
\else
  \newcommand{\PaperTitleText}{A Reproducible Route to a 0.45 Berry--Esseen Bound}
  \newcommand{\PaperRunningTitle}{Toward a 0.45 Berry--Esseen bound}
\fi

\definecolor{BerryBlue}{HTML}{173B57}
\definecolor{BerryLight}{HTML}{EFF5F8}
\definecolor{BerryRule}{HTML}{A7BAC7}
\definecolor{CodeGray}{HTML}{F5F6F7}

\hypersetup{
  hypertexnames=false,
  colorlinks=true,
  linkcolor=BerryBlue,
  citecolor=BerryBlue,
  urlcolor=BerryBlue,
  pdftitle={\PaperTitleText: Regularized Prawitz Estimates and Finite Verification},
  pdfauthor={Kacper Rodziewicz},
  pdfsubject={Berry--Esseen inequality and computer-assisted verification}
}

\pagestyle{fancy}
\fancyhf{}
\fancyhead[LE,LO]{K.\ Rodziewicz}
\fancyhead[RE,RO]{\PaperRunningTitle}
\fancyfoot[C]{\thepage}
\renewcommand{\headrulewidth}{0.35pt}
\fancypagestyle{plain}{
  \fancyhf{}
  \fancyhead[LE,LO]{K.\ Rodziewicz}
  \fancyhead[RE,RO]{\PaperRunningTitle}
  \fancyfoot[C]{\thepage}
  \renewcommand{\headrulewidth}{0.35pt}
}
\fancypagestyle{titlepage}{
  \fancyhf{}
  \fancyfoot[C]{\thepage}
  \renewcommand{\headrulewidth}{0pt}
}

\setlength{\parindent}{1.35em}
\setlength{\parskip}{0pt}
\setlength{\emergencystretch}{2em}
\allowdisplaybreaks
\newcommand{\manualtag}[1]{\tag{#1}}

\newtheorem{theorem}{Theorem}[section]
\newtheorem{lemma}[theorem]{Lemma}
\newtheorem{proposition}[theorem]{Proposition}
\newtheorem{corollary}[theorem]{Corollary}
\theoremstyle{remark}
\newtheorem{remark}[theorem]{Remark}
\newtheorem{reproducibility}[theorem]{Reproducibility record}
\newtheorem*{finaltheorem}{Theorem}
\newtheorem*{finalrecord}{Reproducibility record}

\newcommand{\R}{\mathbb{R}}
\newcommand{\E}{\mathbb{E}}
\newcommand{\Prob}{\mathbb{P}}
\newcommand{\iu}{\mathrm{i}}
\newcommand{\e}{\mathrm{e}}
\newcommand{\dd}{\,\mathrm{d}}
\newcommand{\abs}[1]{\lvert#1\rvert}
\newcommand{\norm}[1]{\left\lVert#1\right\rVert}
\newcommand{\rhozero}{\rho_{0}}
\newcommand{\Cstar}{C^{\ast}}
\newcommand{\Uhat}{\widehat U}
\newcommand{\code}[1]{\texttt{#1}}
\newcommand{\doi}[1]{\href{https://doi.org/#1}{doi:\nolinkurl{#1}}}

\newenvironment{resultbox}
  {\par\medskip\noindent
   \begin{minipage}{\dimexpr\linewidth-2\fboxsep-2\fboxrule\relax}
   \raggedright\color{BerryBlue}\hrule\medskip\color{black}}
  {\medskip\color{BerryBlue}\hrule\color{black}
   \end{minipage}\par\medskip}

\lstset{
  basicstyle=\ttfamily\small,
  backgroundcolor=\color{CodeGray},
  frame=single,
  rulecolor=\color{BerryRule},
  columns=fullflexible,
  keepspaces=true,
  showstringspaces=false,
  breaklines=true,
  xleftmargin=0.5em,
  xrightmargin=0.5em
}

\title{\bfseries\color{BerryBlue}
  \PaperTitleText\\[0.35em]
  \large Regularized Prawitz Estimates and Finite Certificate Verification}
\author{Kacper Rodziewicz}
\date{August 2026}

\begin{document}
\maketitle
\thispagestyle{titlepage}

\begin{abstract}
Let \(X_1,X_2,\ldots\) be independent and identically distributed random
variables with mean zero, variance \(\sigma^2>0\), and finite third absolute
moment.  We develop a complete analytic and computational route to
\[
 \sup_{x\in\R}
 \left|
   \Prob\!\left\{
     \frac{X_1+\cdots+X_n}{\sigma\sqrt n}\le x
   \right\}-\Phi(x)
 \right|
 \le
 0.45\,\frac{\E\abs{X_1}^3}{\sigma^3\sqrt n}.
\]
Known inequalities reduce the problem to
\(2\le n\le8123\) and
\(1\le\rho\le712569/589000\), where
\(\rho=\E\abs{X_1}^3/\sigma^3\).
On this rectangle we combine Prawitz's smoothing inequality with a
Rademacher comparison and two complementary characteristic-function
majorants.  Three analytic devices make rigorous interval evaluation stable:
a positive partial-fraction expansion for a regularized cotangent, a
nonnegative factorization of the Gaussian comparison defect, and an exact
endpoint identity for the Prawitz kernel that gives a quadratic tail.
The resulting criterion is formulated for cell-by-cell verification with
outward-rounded Arb arithmetic.  We specify the certificate language, exact partition
reconstruction, trust boundary, fail-closed checker, and reproducible
execution protocol needed to turn the finite computation into a mathematical
proof.
\ifcertificateaccepted
The archival run reported in the final section completed this verification.
\else
Until the complete public checker reports \code{ACCEPTED\_FULL}, the numerical
conclusion of this typeset version remains conditional.
\fi
\end{abstract}

\begin{center}
\begin{minipage}{0.92\linewidth}
\small
\noindent\textbf{AI assistance disclosure.}
The development of the proof and its computer-assisted verification workflow
was supported by artificial-intelligence assistance.  This manuscript was
also drafted and edited with assistance from OpenAI's ChatGPT using the
GPT-5.6 Sol model.  The named author reviewed the mathematical argument,
computations, and final text and accepts responsibility for the published
claims.
\end{minipage}
\end{center}

\noindent\textbf{2020 Mathematics Subject Classification.}
60F05, 60E10, 65G30.

\noindent\textbf{Keywords.}
Berry--Esseen inequality; Prawitz smoothing inequality; characteristic
functions; interval arithmetic; computer-assisted proof; reproducibility.

\section{Introduction}

For a centered random variable \(Y\) with
\(\E Y^2=1\) and \(\rho=\E\abs{Y}^3<\infty\), write
\[
 S_n=\frac{Y_1+\cdots+Y_n}{\sqrt n},
 \qquad
 \Delta_n=\sup_{x\in\R}\abs{\Prob\{S_n\le x\}-\Phi(x)}.
\]
The classical Berry--Esseen problem asks for the smallest universal constant
\(\Cstar\) such that
\[
 \Delta_n\le \Cstar\,\frac{\rho}{\sqrt n}
\]
for every \(n\) and every admissible distribution.

The best previously published general upper bound is
\(\Cstar\le0.4690\), while Esseen's asymptotic two-point example gives
\[
 \Cstar\ge
 C_E=\frac{\sqrt{10}+3}{6\sqrt{2\pi}}
 =0.409732\ldots.
\]
Thus the value of \(\Cstar\) remains open; see
Shevtsova~\cite{Shevtsova2013} and
Zolotukhin, Nagaev, and Chebotarev~\cite{ZolotukhinNagaevChebotarev2018}
for the preceding bounds and historical context.  The purpose of this paper
is to lower the general upper bound from \(0.4690\) to \(0.45\).

The proof developed here has three layers.  First, published estimates remove
\(n=1\), large third moments, and \(n\ge8124\).  Second, Prawitz's smoothing
inequality is supplied with stable, distribution-free majorants.  Third, the
remaining compact parameter rectangle is covered by finitely many rational
cells whose inequalities are recomputed with directed interval arithmetic.
The logical interface between the last two layers is Proposition~6.3.

The analytic difficulty is concentrated in expressions that are harmless
mathematically but ill-conditioned in floating-point arithmetic.  Near
\(s=0\), subtraction in
\(\cot(\pi s)-1/(\pi s)\) is replaced by a positive series with an explicit
tail.  Near \(u=0\), subtraction between \(\cos^n u\) and its Gaussian
approximation is replaced by a product of nonnegative factors.  Near \(s=1\),
the Prawitz kernel is rewritten in the variable \(h=1-s\), exposing its
linear zero and yielding a contribution of order \(H^2\) on
\([1-H,1]\).

The computer-assisted part is intentionally narrow.  An A100 calculation
supplies parameter advice, but no floating-point GPU value is trusted.
The public checker parses exact rational data, reconstructs every partition,
and recomputes every bound with outward rounding.  It accepts only a complete,
hash-bound bundle for all \(n=2,\ldots,8123\).  Thus search, parallel
scheduling, and checkpoints affect running time, not the mathematical
conclusion.

\section{Reduction to a finite parameter rectangle}

We use three established estimates.  The sharp one-summand result of Bentkus
and Kir\v{s}a~\cite{BentkusKirsa1989} gives
\begin{equation}
 \Delta_1\le 0.3704\,\rho<0.45\,\rho.
 \manualtag{1}\label{eq:one-summand}
\end{equation}
A structural Berry--Esseen estimate due to Shevtsova
\cite{Shevtsova2013} implies
\begin{equation}
 \Delta_n\le
 \frac{0.3322(\rho+0.429)}{\sqrt n}.
 \manualtag{2}\label{eq:structural}
\end{equation}
Finally, an asymptotically optimal estimate with an explicit quadratic
remainder~\cite{Shevtsova2012} gives
\begin{equation}
 \Delta_n\le
 C_E\frac{\rho}{\sqrt n}
 {}+3\left(\frac{\rho}{\sqrt n}\right)^2,
 \qquad
 C_E=\frac{\sqrt{10}+3}{6\sqrt{2\pi}}.
 \manualtag{3}\label{eq:asymptotic}
\end{equation}
By Lyapunov's inequality, \(\rho\ge1\).  Define
\begin{equation}
 \rhozero
 =
 \frac{0.3322\cdot0.429}{0.45-0.3322}
 =
 \frac{712569}{589000}
 =
 1.2097945670\ldots.
 \manualtag{4}\label{eq:rhozero}
\end{equation}

\begin{lemma}[Large third moments]\label{lem:large-rho}
If \(\rho\ge\rhozero\), then
\[
 \Delta_n\le0.45\,\frac{\rho}{\sqrt n}
\]
for every \(n\ge1\).
\end{lemma}

\begin{proof}
By \eqref{eq:structural}, it is enough to check
\(0.3322(\rho+0.429)\le0.45\rho\), which is equivalent to
\(\rho\ge\rhozero\).
\end{proof}

\begin{lemma}[Large sample sizes]\label{lem:large-n}
If \(1\le\rho\le\rhozero\) and \(n\ge8124\), then
\[
 \Delta_n\le0.45\,\frac{\rho}{\sqrt n}.
\]
\end{lemma}

\begin{proof}
Equation~\eqref{eq:asymptotic} gives
\[
 \frac{\sqrt n\,\Delta_n}{\rho}
 \le C_E+\frac{3\rho}{\sqrt n}
 \le C_E+\frac{3\rhozero}{\sqrt n}.
\]
The right-hand side decreases with \(n\).  The strict inequality at
\(n=8124\) follows from the exact rational calculation in
Appendix~\ref{app:threshold}.
\end{proof}

\begin{corollary}[Finite rectangle]\label{cor:finite-rectangle}
To prove \(\Cstar\le0.45\), it is enough to treat
\begin{equation}
 2\le n\le8123,
 \qquad
 1\le\rho\le\rhozero.
 \manualtag{5}\label{eq:finite-rectangle}
\end{equation}
\end{corollary}

\begin{proof}
Use \eqref{eq:one-summand} for \(n=1\),
Lemma~\ref{lem:large-rho} for \(\rho\ge\rhozero\), and
Lemma~\ref{lem:large-n} for \(n\ge8124\) in the remaining range.
\end{proof}

\section{The Prawitz functional}

Let
\[
 f_n(t)=\E\e^{\iu tS_n},
 \qquad
 r_n(t)=\abs{f_n(t)-\e^{-t^2/2}}.
\]
For \(0<s<1\), define the Prawitz kernel
\begin{equation}
 K(s)=
 \frac{1-s}{2}
 {}+\frac{\iu}{2}
 \left((1-s)\cot(\pi s)+\frac1\pi\right),
 \manualtag{6}\label{eq:kernel}
\end{equation}
and extend it continuously by \(K(1)=0\).
Prawitz's smoothing inequality~\cite{Prawitz1972,KorolevShevtsova2012}
states that, for \(n>1\), \(T>0\), and \(0<t_0\le1\),
\begin{equation}
\begin{aligned}
 \Delta_n\le{}&
 2\int_0^{t_0}\abs{K(s)}r_n(Ts)\dd s\\
 &{}+2\int_{t_0}^{1}\abs{K(s)}\abs{f_n(Ts)}\dd s\\
 &{}+2\int_0^{t_0}
 \left|K(s)-\frac{\iu}{2\pi s}\right|
 \e^{-T^2s^2/2}\dd s
 {}+\frac1\pi\int_{t_0}^{\infty}
 \frac{\e^{-T^2s^2/2}}{s}\dd s.
\end{aligned}
 \manualtag{7}\label{eq:prawitz}
\end{equation}
The final term equals
\[
 \frac{1}{2\pi}E_1\!\left(\frac{T^2t_0^2}{2}\right),
 \qquad
 E_1(x)=\int_x^\infty\frac{\e^{-y}}{y}\dd y.
\]

Set
\begin{equation}
 B(\rho)=
 \sqrt{\frac{\rho^2}{2}
       +\frac{\rho\sqrt{\rho^2+8}}{2}-2}.
 \manualtag{8}\label{eq:B}
\end{equation}
Write \(\varphi(u)=\E\e^{\iu uY_1}\), so that
\(f_n(t)=\varphi(t/\sqrt n)^n\).  The comparison theorem of Mattner and
Shevtsova~\cite{MattnerShevtsova2019} gives the one-summand estimate
\[
 \abs{\varphi(u)-\cos u}
 \le \frac{B(\rho)\abs u^3}{6}.
\]
Put
\[
 u=\frac{t}{\sqrt n},
 \qquad
 c=\abs{\cos u},
 \qquad
 \delta=\frac{B(\rho)\abs u^3}{6}.
\]
The power factorization, together with
\(\abs{\varphi(u)}\le1\) and \(c\le1\), first yields
\begin{equation}
 \left|
 f_n(t)-\cos^n\!\left(\frac{t}{\sqrt n}\right)
 \right|
 \le
 n\delta
 =\frac{B(\rho)\abs t^3}{6\sqrt n}.
 \manualtag{9}\label{eq:rademacher-comparison}
\end{equation}
Moreover,
\(\abs{\varphi(u)}\le\min\{1,c+\delta\}=:a\).  Define
\[
 D_{\mathrm{pow}}(n,\rho,t)
 =
 \delta\sum_{j=0}^{n-1}a^jc^{\,n-1-j},
 \qquad
 D(n,\rho,t)
 =
 \min\{2,n\delta,D_{\mathrm{pow}}(n,\rho,t)\}.
\]
The factorization of \(z^n-w^n\), with
\(\abs{\varphi(u)}\le a\), gives the bound \(D_{\mathrm{pow}}\).
The triangle inequality gives the bound \(2\), while
\eqref{eq:rademacher-comparison} gives \(n\delta\).  Taking the minimum gives
\begin{equation}
 \left|
 f_n(t)-\cos^n\!\left(\frac{t}{\sqrt n}\right)
 \right|
 \le D(n,\rho,t).
 \manualtag{10}\label{eq:power-bound}
\end{equation}

Let \(\theta_0\) be the unique root in \((\pi,2\pi)\) of
\[
 \theta^2+2\theta\sin\theta+6(\cos\theta-1)=0,
\]
and define
\[
 \kappa=
 \frac{\cos\theta_0-1+\theta_0^2/2}{\theta_0^3}.
\]
For \(z\ge0\), set
\begin{equation}
 \psi(z)=
 \begin{cases}
  z^2/2-\kappa z^3, & 0\le z\le\theta_0,\\
  1-\cos z,         & \theta_0<z\le2\pi,\\
  0,                & z>2\pi.
 \end{cases}
 \manualtag{11}\label{eq:psi}
\end{equation}

\begin{lemma}[Bounds for \(\psi\)]\label{lem:psi-bounds}
The function \(\psi\) in \eqref{eq:psi} is continuous on
\([0,\infty)\) and satisfies
\[
 0\le\psi(z)\le2
 \qquad(z\ge0).
\]
Consequently, for \(\rho\ge1\), the base in \eqref{eq:MP} belongs
to \([0,1]\).
\end{lemma}

\begin{proof}
Write
\[
 F(z)=z^2+2z\sin z+6(\cos z-1)
\]
and, for \(z>0\),
\[
 r(z)=\frac{\cos z-1+z^2/2}{z^3}.
\]
Direct differentiation gives
\[
 r'(z)=-\frac{F(z)}{2z^4},
 \qquad
 F'(z)=4\cos(z/2)
 \bigl(z\cos(z/2)-2\sin(z/2)\bigr).
\]
For \(0<z<\pi\), the inequality \(\tan(z/2)>z/2\) implies
\(F'(z)<0\), and hence \(F(z)<0\).
For \(\pi<z<2\pi\), both factors in the displayed product for
\(F'\) are negative, so \(F'(z)>0\).  Moreover,
\[
 F(\pi)=\pi^2-12<0,
 \qquad
 F(2\pi)=4\pi^2>0.
\]
The intermediate value theorem therefore gives a zero in \((\pi,2\pi)\),
and strict monotonicity there makes it unique.  It follows that
\(F(z)<0\) for \(0<z<\theta_0\).  Thus \(r\) is increasing there
and
\[
 r(z)\le r(\theta_0)=\kappa
 \qquad(0<z\le\theta_0).
\]
Therefore
\[
 z^2/2-\kappa z^3\le1-\cos z\le2.
\]
Moreover, \(1-\cos\theta_0<\theta_0^2/2\) gives \(\kappa>0\), and
\[
 \kappa\theta_0
 =
 \frac12-\frac{1-\cos\theta_0}{\theta_0^2}
 \le\frac12.
\]
Hence
\[
 z^2/2-\kappa z^3
 =
 z^2(1/2-\kappa z)\ge0
 \qquad(0\le z\le\theta_0).
\]
The bounds on the other branches are immediate.
The defining identity for \(\kappa\) gives continuity at
\(\theta_0\), and both adjacent branches vanish at \(2\pi\).
Finally, \((\rho+1)^2\ge4\).
\end{proof}

For \(\epsilon>0\), write
\[
 \Psi(t,\epsilon)=\frac{\psi(\epsilon\abs t)}{\epsilon^2}.
\]
For \(\delta=1\), Theorem~2.2 of Shevtsova~\cite{Shevtsova2012}
states
\[
 \abs{f_n(t)}
 \le
 \left(1-\frac2n\Psi(t,\epsilon_n)\right)^{n/2},
 \qquad
 \epsilon_n=\frac{\rho+\E\abs Y}{\sqrt n}.
\]
The cited result also proves that \(\Psi(t,\epsilon)\) is nonincreasing
in \(\epsilon\) for fixed \(t\).  Since Cauchy--Schwarz gives
\(\E\abs Y\le1\), we may replace \(\epsilon_n\) by the larger value
\((\rho+1)/\sqrt n\).  After the scaling
\begin{equation}
 T=\frac{q\sqrt n}{\rho+1},
 \manualtag{12}\label{eq:T-scaling}
\end{equation}
this gives
\begin{equation}
 \abs{f_n(Ts)}
 \le
 M_P(n,\rho,Ts)
 :=
 \left(
 1-\frac{2\psi(qs)}{(\rho+1)^2}
 \right)^{n/2}.
 \manualtag{13}\label{eq:MP}
\end{equation}
Indeed, \(((\rho+1)/\sqrt n)Ts=qs\), and hence
\(n^{-1}\Psi(Ts,(\rho+1)/\sqrt n)
=\psi(qs)/(\rho+1)^2\), exactly the base in \eqref{eq:MP}.
Lemma~\ref{lem:psi-bounds} makes the real half-power in
\eqref{eq:MP} well defined.

Define
\begin{equation}
 R_{n,\rho}(t)
 =
 \min\left\{
 2,\,
 \left|
 \cos^n\!\left(\frac{t}{\sqrt n}\right)-\e^{-t^2/2}
 \right|
 +D(n,\rho,t)
 \right\},
 \manualtag{14}\label{eq:R}
\end{equation}
\begin{equation}
 M_R(n,\rho,t)
 =
 \min\left\{
 1,\,
 \left|\cos\!\left(\frac{t}{\sqrt n}\right)\right|^n
 +D(n,\rho,t)
 \right\},
 \manualtag{15}\label{eq:MR}
\end{equation}
and
\begin{equation}
 M_{n,\rho}(t)
 =
 \min\{M_R(n,\rho,t),M_P(n,\rho,t)\}.
 \manualtag{16}\label{eq:M}
\end{equation}
Then
\begin{equation}
 r_n(t)\le R_{n,\rho}(t),
 \qquad
 \abs{f_n(t)}\le M_{n,\rho}(t)\le1.
 \manualtag{17}\label{eq:majorants}
\end{equation}

Substitution into \eqref{eq:prawitz} gives
\begin{equation}
\begin{aligned}
 U_n(\rho;T,t_0)
 ={}&
 2\int_0^{t_0}\abs{K(s)}R_{n,\rho}(Ts)\dd s\\
 &{}+2\int_{t_0}^{1}\abs{K(s)}M_{n,\rho}(Ts)\dd s\\
 &{}+2\int_0^{t_0}\abs{L(s)}\e^{-T^2s^2/2}\dd s
 {}+\frac1{2\pi}E_1\!\left(\frac{T^2t_0^2}{2}\right),
\end{aligned}
 \manualtag{18}\label{eq:U}
\end{equation}
where \(L(s)=K(s)-\iu/(2\pi s)\).  For every admissible
distribution,
\begin{equation}
 \Delta_n\le U_n(\rho;T,t_0).
 \manualtag{19}\label{eq:universal-U}
\end{equation}

\section{Stable analytic majorants}

\subsection{The coefficient \texorpdfstring{\(B(\rho)\)}{B(rho)}}

The direct formula \eqref{eq:B} contains cancellation at \(\rho=1\).
The following factorization removes it.

\begin{lemma}\label{lem:stable-B}
For every \(\rho\ge1\),
\begin{equation}
 2B(\rho)^2
 =
 (\rho^2-1)
 \left(
 1+\frac{\rho^2+9}{\rho\sqrt{\rho^2+8}+3}
 \right).
 \manualtag{20}\label{eq:stable-B}
\end{equation}
Moreover, \(B(\rho)\ge0\) and \(B\) is increasing on
\([1,\infty)\).
\end{lemma}

\begin{proof}
Rationalizing the radical in \eqref{eq:B} gives
\eqref{eq:stable-B}.  For \(\rho>1\),
\[
 \frac{\dd}{\dd\rho}B(\rho)^2
 =
 \rho+\frac{\rho^2+4}{\sqrt{\rho^2+8}}>0.
\]
Since \(B\) is the nonnegative square root in \eqref{eq:B} and
\(B(1)=0\), monotonicity follows by continuity.
\end{proof}

\subsection{A positive recursion for power differences}

\begin{lemma}\label{lem:positive-recursion}
For \(a,c\ge0\) and every integer \(m\ge1\), let
\[
 S_n(a,c)=\sum_{j=0}^{n-1}a^jc^{\,n-1-j}.
\]
Then
\[
 S_{2m}=S_m(a,c)(a^m+c^m),
 \qquad
 S_{2m+1}=a^{2m}+cS_{2m}(a,c).
\]
\end{lemma}

\begin{proof}
Split the first sum at \(m\), and separate the term \(a^{2m}\) in
the second.  The recursion uses only addition and multiplication of
nonnegative quantities and requires no division by \(a-c\); it is therefore
well suited to outward-rounded evaluation.
\end{proof}

\subsection{The regularized cotangent}

For \(0<s<1\), define
\begin{equation}
 d(s)=\frac1{\pi s}-\cot(\pi s).
 \manualtag{21}\label{eq:d}
\end{equation}

\begin{lemma}[Positive partial fractions]\label{lem:cotangent}
For \(0<x<\pi\),
\begin{equation}
 \frac1x-\cot x
 =
 2x\sum_{k=1}^{\infty}
 \frac1{k^2\pi^2-x^2}.
 \manualtag{22}\label{eq:cot-series}
\end{equation}
Every summand is positive and strictly increasing in \(x\).  For each
integer \(m\ge1\),
\begin{equation}
 \frac1x-\cot x
 \le
 2x\sum_{k=1}^{m}
 \frac1{k^2\pi^2-x^2}
 +
 \frac{2x}{
   m\bigl(\pi^2-x^2/(m+1)^2\bigr)}.
 \manualtag{23}\label{eq:cot-tail}
\end{equation}
Consequently, \(d\) is nonnegative and strictly increasing on \((0,1)\).
\end{lemma}

\begin{proof}
The logarithmic derivative of Euler's product
\[
 \frac{\sin x}{x}
 =
 \prod_{k=1}^{\infty}
 \left(1-\frac{x^2}{k^2\pi^2}\right),
\]
with termwise differentiation justified uniformly on compact subsets of
\((-\pi,\pi)\), gives \eqref{eq:cot-series}.  For fixed \(k\),
\[
 \frac{\dd}{\dd x}
 \frac{x}{k^2\pi^2-x^2}
 =
 \frac{k^2\pi^2+x^2}{(k^2\pi^2-x^2)^2}>0.
\]
For \(k\ge m+1\),
\[
 \frac1{k^2\pi^2-x^2}
 \le
 \frac1{k^2\bigl(\pi^2-x^2/(m+1)^2\bigr)}.
\]
Using
\(\sum_{k=m+1}^{\infty}k^{-2}\le\int_m^\infty t^{-2}\dd t=1/m\)
proves \eqref{eq:cot-tail}.  Substitution \(x=\pi s\) proves the
final assertion.
\end{proof}

The regularized kernel has the exact form
\begin{equation}
 L(s)
 =
 \frac{1-s}{2}\bigl(1-\iu d(s)\bigr),
 \qquad
 2\abs{L(s)}
 =
 (1-s)\sqrt{1+d(s)^2}.
 \manualtag{24}\label{eq:L}
\end{equation}

\begin{corollary}[Uniform \(L\)-bound]\label{cor:L-bound}
If \(0<s_-\le s_+<1\), then
\begin{equation}
 \sup_{s\in[s_-,s_+]}2\abs{L(s)}
 \le
 (1-s_-)\sqrt{1+d(s_+)^2}.
 \manualtag{25}\label{eq:L-bound}
\end{equation}
The value \(d(s_+)\) may be replaced by the finite upper bound
in \eqref{eq:cot-tail}.
\end{corollary}

\begin{proof}
Use \eqref{eq:L}, Lemma~\ref{lem:cotangent}, and monotonicity of the
square root.
\end{proof}

\subsection{A stable Gaussian comparison}

Directly subtracting \(\cos^n u\) and \(\e^{-nu^2/2}\) is poorly
conditioned near zero.  The next identity separates the small defect from a
Gaussian factor.

\begin{lemma}[Gaussian defect]\label{lem:gaussian-defect}
Let \(n\ge1\) be an integer.  For \(0\le u<\pi/2\), define
\begin{equation}
 h(u)=-\log(\cos u)-\frac{u^2}{2}.
 \manualtag{26}\label{eq:h}
\end{equation}
Then \(h(0)=0\), \(h(u)\ge0\), and \(h\) is increasing.  Moreover,
\begin{equation}
 \left|\cos^n u-\e^{-nu^2/2}\right|
 =
 \e^{-nu^2/2}\left(1-\e^{-nh(u)}\right).
 \manualtag{27}\label{eq:gaussian-defect}
\end{equation}
\end{lemma}

\begin{proof}
Since \(h'(u)=\tan u-u\) and
\((\tan u-u)'=\tan^2u\ge0\), one has \(h'\ge0\) and \(h\ge0\).
Positivity of \(\cos u\) permits taking the logarithm and gives
\[
 \cos^n u=\e^{-nu^2/2}\e^{-nh(u)}.
\]
Equation~\eqref{eq:gaussian-defect} follows.
\end{proof}

\begin{corollary}[Rectangular-cell bound]\label{cor:gaussian-cell}
Let \(n\ge1\) be an integer,
\(\rho\in[\rho_-,\rho_+]\), \(s\in[s_-,s_+]\), and \(q>0\).  Put
\[
 u_-=\frac{qs_-}{\rho_++1},
 \qquad
 u_+=\frac{qs_+}{\rho_-+1}.
\]
If \(0\le u_-\le u_+<\pi/2\), then throughout the cell
\begin{equation}
\begin{aligned}
 \left|
 \cos^n\!\left(\frac{qs}{\rho+1}\right)
 -
 \exp\!\left(
 -\frac{nq^2s^2}{2(\rho+1)^2}
 \right)
 \right|
 &\le
 \e^{-nu_-^2/2}\left(1-\e^{-nh(u_+)}\right)\\
 &\le
 \e^{-nu_-^2/2}\min\{1,nh(u_+)\}.
\end{aligned}
 \manualtag{28}\label{eq:gaussian-cell}
\end{equation}
\end{corollary}

\begin{proof}
Throughout the rectangle,
\[
 u_-\le\frac{qs}{\rho+1}\le u_+.
\]
In \eqref{eq:gaussian-defect}, the Gaussian factor decreases with \(u\),
while \(1-\e^{-nh(u)}\) increases.  Bound the two factors at opposite
endpoints.  The final inequality uses
\(1-\e^{-x}\le\min\{1,x\}\) for \(x\ge0\).
\end{proof}

\section{Analytic control of the endpoint \texorpdfstring{\(s=1\)}{s=1}}

The second integral in \eqref{eq:U} reaches \(s=1\).  A rectangular
enclosure that forgets \(K(1)=0\) introduces a first-order artifact.
The endpoint can instead be handled exactly.

\begin{lemma}[Endpoint identity]\label{lem:endpoint-identity}
For \(0<h<1\),
\begin{equation}
 K(1-h)=\frac h2+\frac{\iu h}{2}d(h),
 \qquad
 2\abs{K(1-h)}=h\sqrt{1+d(h)^2}.
 \manualtag{29}\label{eq:endpoint-identity}
\end{equation}
In particular, \(K(1)=0\) by continuity.
\end{lemma}

\begin{proof}
Using \(\cot(\pi-\pi h)=-\cot(\pi h)\) in \eqref{eq:kernel} gives
\[
 K(1-h)
 =
 \frac h2+
 \frac{\iu}{2}
 \left(-h\cot(\pi h)+\frac1\pi\right)
 =
 \frac h2+\frac{\iu h}{2}d(h).
\]
Taking the modulus proves \eqref{eq:endpoint-identity}.
Equation~\eqref{eq:cot-series} gives \(d(h)=O(h)\) as \(h\downarrow0\),
and hence \(K(1-h)\to0\).
\end{proof}

\begin{lemma}[Uniform endpoint shape]\label{lem:endpoint-shape}
If \(0\le h\le H\le1/2\), then
\begin{equation}
 2\abs{K(1-h)}
 \le
 h\sqrt{1+d(H)^2}.
 \manualtag{30}\label{eq:endpoint-shape}
\end{equation}
\end{lemma}

\begin{proof}
For \(h=0\), the assertion follows from \(K(1)=0\).
If \(0<h\le H\), Lemma~\ref{lem:cotangent} gives
\(0\le d(h)\le d(H)\); now apply \eqref{eq:endpoint-identity}.
\end{proof}

\begin{theorem}[Quadratic endpoint tail]\label{thm:quadratic-tail}
Let \(n\ge1\) be an integer, let \(\rho\ge1\), \(T>0\), and let
\(0<H\le1/2\), \(s_\ast=1-H\).  Then
\begin{equation}
 \int_{s_\ast}^{1}
 2\abs{K(s)}M_{n,\rho}(Ts)\dd s
 \le
 \frac{H^2}{2}\sqrt{1+d(H)^2}.
 \manualtag{31}\label{eq:quadratic-tail}
\end{equation}
\end{theorem}

\begin{proof}
By \eqref{eq:majorants}, \(0\le M_{n,\rho}\le1\).  Substitute
\(s=1-h\), apply Lemma~\ref{lem:endpoint-shape}, and integrate:
\[
 \int_{s_\ast}^{1}
 2\abs{K(s)}M_{n,\rho}(Ts)\dd s
 \le
 \int_0^H h\sqrt{1+d(H)^2}\dd h.
\]
\end{proof}

\begin{remark}[Sharp order]\label{rem:sharp-order}
Equation~\eqref{eq:cot-series} gives
\[
 d(H)=\frac{\pi H}{3}+O(H^3),
 \qquad
 \frac{H^2}{2}\sqrt{1+d(H)^2}
 =
 \frac{H^2}{2}+O(H^4).
\]
A bound that treats the endpoint imaginary part as a nonzero constant would
produce an artificial \(O(H)\) contribution.  The quadratic estimate restores
the correct local order of the kernel.
\end{remark}

\begin{lemma}[Choice of cutoff]\label{lem:cutoff}
Let \(r\ge1\) be an integer and let
\[
 \max\{t_0,1/2\}\le s_1,\ldots,s_r<1.
\]
For each \(j\), suppose that \(P_j\) is a rigorous upper bound for
\[
 \int_{t_0}^{s_j}
 2\abs{K(s)}M_{n,\rho}(Ts)\dd s.
\]
Then
\begin{equation}
 \int_{t_0}^{1}
 2\abs{K(s)}M_{n,\rho}(Ts)\dd s
 \le
 \min_{1\le j\le r}
 \left\{
 P_j+
 \frac{(1-s_j)^2}{2}
 \sqrt{1+d(1-s_j)^2}
 \right\}.
 \manualtag{32}\label{eq:cutoff}
\end{equation}
\end{lemma}

\begin{proof}
For each \(j\), split the integral at \(s_j\).
Since \(H_j=1-s_j\in(0,1/2]\), Theorem~\ref{thm:quadratic-tail}
applies to the second part.  Every expression inside the minimum is therefore
an upper bound for the complete integral, and the finite minimum is valid.
\end{proof}

\begin{remark}[Verifier cutoffs]\label{rem:verifier-cutoffs}
The verifier uses the exact dyadic cutoffs
\(s_j=1-2^{-j}\) for \(j=1,\ldots,14\), retaining only eligible
cutoffs.  It also tests the left endpoint of the final certificate
subinterval when that endpoint is at least \(1/2\).  The value of \(d\) is
bounded by \eqref{eq:cot-tail} with \(m=100\).
\end{remark}

\section{A finite verification criterion}

Fix \(n\in\{2,\ldots,8123\}\) and
\(I=[\rho_-,\rho_+]\subset[1,\rhozero]\).
Choose positive rational \(q,t_0,\eta,\lambda\) such that
\[
 0<\lambda\le\min\{t_0,1/2\},
 \qquad t_0<1,
 \qquad
 0<\eta\le\min\{t_0,1/2\},
 \qquad
 \frac{q\eta}{\rho_-+1}\le1,
\]
and use \(T(\rho)=q\sqrt n/(\rho+1)\) as in
\eqref{eq:T-scaling}.  For every \(\rho\in I\),
\[
 \frac{T(\rho)\eta}{\sqrt n}
 =
 \frac{q\eta}{\rho+1}
 \le
 \frac{q\eta}{\rho_-+1}
 \le1.
\]

\begin{lemma}[First integral near zero]\label{lem:first-zero}
Let \(n\ge1\) be an integer, \(\rho\ge1\), \(T>0\), and
\(0<\eta\le1/2\).  If \(T\eta/\sqrt n\le1\), then
\begin{equation}
 2\int_0^\eta
 \abs{K(s)}R_{n,\rho}(Ts)\dd s
 \le
 \frac{B(\rho)T^3\eta^3}{9\sqrt n}
 +
 \frac{T^4\eta^4}{12n}.
 \manualtag{33}\label{eq:first-zero}
\end{equation}
\end{lemma}

\begin{proof}
First justify the kernel estimate.  If \(0<s<1/2\), then
\(\tan(\pi s)\ge\pi s\), so
\(0\le\cot(\pi s)\le1/(\pi s)\).  Equation~\eqref{eq:kernel}
gives
\[
 \abs{\Re K(s)}\le\frac12,
 \qquad
 \abs{\Im K(s)}\le\frac1{2\pi s}.
\]
Thus
\[
 \abs{K(s)}
 \le\abs{\Re K(s)}+\abs{\Im K(s)}
 \le\frac12+\frac1{2\pi s}\le\frac1s.
\]
At \(s=1/2\), the same estimate follows directly from
\(\cot(\pi/2)=0\); the endpoint \(s=0\) is understood by the displayed
majorant.

Put \(u=Ts/\sqrt n\).  On \(0\le s\le\eta\), one has
\(0\le u\le1\).  Lemma~\ref{lem:gaussian-defect} gives
\(\cos u\le\e^{-u^2/2}\).  Moreover,
\[
 \cos u\ge1-\frac{u^2}{2},
 \qquad
 \e^{-u^2/2}
 \le
 1-\frac{u^2}{2}+\frac{u^4}{8},
\]
and hence
\[
 0\le\e^{-u^2/2}-\cos u\le\frac{u^4}{8}.
\]
Both factors lie in \([0,1]\).  The power factorization therefore gives
\[
 \left|\e^{-nu^2/2}-\cos^n u\right|
 \le
 n\left(\e^{-u^2/2}-\cos u\right)
 \le\frac{nu^4}{8}\le\frac{nu^4}{6}.
\]
From the definition preceding \eqref{eq:power-bound},
\[
 D(n,\rho,Ts)\le n\delta
 =
 \frac{B(\rho)T^3s^3}{6\sqrt n}.
\]
Equation~\eqref{eq:R} now gives
\[
 R_{n,\rho}(Ts)
 \le
 \frac{B(\rho)T^3s^3}{6\sqrt n}
 +
 \frac{T^4s^4}{6n}.
\]
Multiply by \(2\abs{K(s)}\le2/s\) and integrate.
\end{proof}

\begin{lemma}[Third integral near zero]\label{lem:third-zero}
For \(0<\lambda\le\min\{t_0,1/2\}\),
\begin{equation}
 2\int_0^\lambda
 \abs{L(s)}\e^{-T^2s^2/2}\dd s
 \le
 \lambda\sqrt{1+\frac4{\pi^2}}.
 \manualtag{34}\label{eq:third-zero}
\end{equation}
\end{lemma}

\begin{proof}
For \(0<s\le\lambda\le1/2\), Lemma~\ref{lem:cotangent} gives
\(d(s)\le d(1/2)=2/\pi\).  Equation~\eqref{eq:L} implies
\[
 2\abs{L(s)}
 =
 (1-s)\sqrt{1+d(s)^2}
 \le
 \sqrt{1+4/\pi^2}.
\]
Since \(\e^{-T^2s^2/2}\le1\), integration proves
\eqref{eq:third-zero}; the endpoint \(s=0\) follows by continuity.
\end{proof}

\subsection{Uniform evaluation on a cell}

Partition \([\eta,t_0]\), \([t_0,1]\), and
\([\lambda,t_0]\) into finitely many rational subintervals.
On every proper subrectangle \(I\times[s_-,s_+]\), all quantities are
enclosed by outward-rounded interval arithmetic.

Lemma~\ref{lem:stable-B} bounds \(B(\rho)\) by \(B(\rho_+)\).
The scaled variable satisfies
\[
 \frac{qs_-}{\rho_++1}
 \le
 \frac{qs}{\rho+1}
 \le
 \frac{qs_+}{\rho_-+1}.
\]
Equations~\eqref{eq:B}--\eqref{eq:majorants},
Lemma~\ref{lem:positive-recursion}, and
Corollary~\ref{cor:gaussian-cell} give uniform upper bounds for
\(R_{n,\rho}\) and \(M_{n,\rho}\).  If the upper endpoint of the
scaled interval is not below \(\pi/2\), the difference in \eqref{eq:R}
is evaluated directly by interval arithmetic.  Definition~\eqref{eq:psi}
is evaluated on cells contained in one branch; a subinterval crossing
\(\theta_0\) or \(2\pi\) uses the global lower bound
\(\psi\ge0\) from Lemma~\ref{lem:psi-bounds}.  Continuity of \(\psi\)
at both branch points and continuity of the other elementary expressions
show that the proper integrands are continuous.

Equation~\eqref{eq:kernel} bounds the kernel in the first two proper
integrals.  Equations~\eqref{eq:cot-tail}--\eqref{eq:L-bound} bound the
regularized kernel in the third.  Each proper integral is enclosed by a
finite upper sum.  For the second integral, every prefix sum is combined
with Theorem~\ref{thm:quadratic-tail} and Lemma~\ref{lem:cutoff}.

The near-zero part of the first integral is bounded uniformly by
\eqref{eq:first-zero}, using \(B(\rho_+)\) and
\[
 T_{\max}=\frac{q\sqrt n}{\rho_-+1}.
\]
The third near-zero part is bounded by \eqref{eq:third-zero}.
Finally,
\[
 E_1'(x)=-\frac{\e^{-x}}x<0.
\]
Since \(T(\rho)\) is smallest at \(\rho=\rho_+\), the last term in
\eqref{eq:U} is bounded uniformly by
\[
 \frac1{2\pi}
 E_1\!\left(
  \frac{nq^2t_0^2}{2(\rho_++1)^2}
 \right).
\]
Let \(\Uhat_{n,I}\) denote the sum of these six certified
contributions.  Then
\[
 U_n(\rho;T(\rho),t_0)\le\Uhat_{n,I}
 \qquad(\rho\in I).
\]

The implementation isolates the two branch points with the rational guards
\[
 \frac{3995895679}{10^9}<\theta_0<
 \frac{3995895680}{10^9},
 \qquad
 \frac{6283185307}{10^9}<2\pi<
 \frac{6283185308}{10^9}.
\]
These inequalities are themselves certified at initialization.

\begin{proposition}[Cell criterion]\label{prop:cell-criterion}
If
\begin{equation}
 \frac{\sqrt n}{\rho_-}\,\Uhat_{n,I}
 \le\frac9{20},
 \manualtag{35}\label{eq:cell-criterion}
\end{equation}
then
\[
 \frac{\sqrt n\,\Delta_n}{\rho}\le0.45
\]
for every \(\rho\in I\) and every admissible distribution with third
absolute moment \(\rho\).
\end{proposition}

\begin{proof}
By \eqref{eq:universal-U},
\(\Delta_n\le\Uhat_{n,I}\) throughout the cell.  Since
\(\rho\ge\rho_-\),
\[
 \frac{\sqrt n\,\Delta_n}{\rho}
 \le
 \frac{\sqrt n\,\Uhat_{n,I}}{\rho_-}
 \le\frac9{20}=0.45.
\]
\end{proof}

\begin{theorem}[Finite completion criterion]\label{thm:finite-completion}
Suppose that, for every \(n=2,\ldots,8123\), the interval
\([1,\rhozero]\) is covered by finitely many cells satisfying
\eqref{eq:cell-criterion}.  Then \(\Cstar\le0.45\).
\end{theorem}

\begin{proof}
Apply Proposition~\ref{prop:cell-criterion} on the finite rectangle
\eqref{eq:finite-rectangle}, and use
Corollary~\ref{cor:finite-rectangle} outside it.
\end{proof}

\section{Computer-assisted verification}
\label{sec:computer-verification}

The finite covering in Theorem~\ref{thm:finite-completion} is implemented as
a certificate-checking computation.  The proof object is separated from the
procedure used to discover it: GPU searches and heuristic parameter selection
may propose a certificate, but acceptance depends on a compact checker that
reconstructs every partition and recomputes every bound with outward-rounded
interval arithmetic.

\subsection{Scope and proof objects}

For each integer \(n=2,\ldots,8123\), a certificate shard contains an
ordered partition
\[
 1=\rho_{n,0}<\rho_{n,1}<\cdots<\rho_{n,m_n}=\rhozero.
\]
Each cell
\(I_{n,j}=[\rho_{n,j-1},\rho_{n,j}]\) carries exact rational
parameters \(q,t_0,\eta,\lambda\) and a compact grid recipe with two
integers, \code{pieces} and \code{power}.  The checker requires literal
equality of consecutive rational endpoints.  Gaps, overlaps, reversed
intervals, and degenerate cells are rejected.

For each cell the checker reconstructs all six contributions to
\(\Uhat_{n,I}\): the two near-zero estimates
\eqref{eq:first-zero} and \eqref{eq:third-zero}, three proper upper sums,
and the \(E_1\) tail.  It accepts the cell only if the outward-rounded upper
endpoint of the left-hand side of \eqref{eq:cell-criterion} does not exceed
the lower endpoint of the exact rational \(9/20\).

\subsection{Exact certificate language}

A shard has format \code{berry-esseen-045-cert-v2}.  All real-valued
certificate data are canonical rational strings: an integer or a reduced
fraction \(a/b\) with \(b>0\).  JSON floating-point numbers, duplicate keys,
unknown fields, noncanonical fractions, and nonstandard constants such as
\code{NaN} or \code{Infinity} are rejected.  A cell contains exactly
\[
 \code{rho},\ \code{q},\ \code{t0},\ \code{eta},\
 \code{lambda},\ \code{grid},
\]
where \code{rho} contains two endpoints and
\(\code{grid}=\{\code{pieces},\code{power}\}\), with
\[
 1\le\code{pieces}\le4096,
 \qquad
 1\le\code{power}\le8.
\]
The analytic guards
\[
 0<\lambda\le\min\{t_0,1/2\},\qquad t_0<1,\qquad
 0<\eta\le\min\{t_0,1/2\},\qquad
 q\eta\le\rho_-+1
\]
are checked again from the parsed rationals.

\subsection{Deterministic reconstruction of the grids}

No floating-point grid nodes occur in the proof object.  Let \([a,b]\) be
one of the proper integration intervals, \(p=\code{pieces}\), and
\(d=\code{power}\).  A concentrated grid is the sorted union of the
following exact rational points:

\begin{enumerate}[label=(\roman*),leftmargin=2.2em]
\item the polynomial mesh
\[
 a+(b-a)\frac{i^d}{M^d},
 \qquad
 M=\max\{8,\lfloor p/2\rfloor\},
 \qquad
 0\le i\le M;
\]
\item the uniform safety mesh with
\(\max\{4,\lfloor p/8\rfloor\}\) subintervals; and
\item for \(n\ge16\), a mesh in the scaled variable \(Ts\).
Writing
\[
 \underline{\sqrt n}
 =
 \frac{\lfloor10^6\sqrt n\rfloor}{10^6},
 \qquad
 D=\max\{8,\lfloor p/4\rfloor\},
\]
where the floor is computed by integer square root, the inserted points are
\[
 \frac{\rho_{\mathrm{mid}}+1}
      {q\,\underline{\sqrt n}}\frac{k}{D},
 \qquad
 0\le k\le10D,
\]
whenever they lie strictly inside \((a,b)\).
\end{enumerate}

For \(n<16\) and \(d=1\), the concentrated grid reduces to the exact
uniform \(p\)-partition.  The first and third proper integrals use
concentrated grids; the second uses the exact uniform \(p\)-partition.
Every grid is augmented with the four rational guards for
\(\theta_0/q\) and \(2\pi/q\) displayed in Section~6.  Thus branch crossings
of \(\psi(qs)\) are isolated deterministically.  The second grid is also
augmented by the dyadic endpoint cutoffs described in
Remark~\ref{rem:verifier-cutoffs}.

\subsection{Outward-rounded evaluation}

The trusted numerical core uses Arb through
\code{python-flint==0.9.0}.  Exact rational points are passed through FLINT
rationals; an interval \([a,b]\) is represented by the ball with exact
midpoint \((a+b)/2\) and radius \((b-a)/2\).  The publication configuration
uses \(192\) bits of precision.  All elementary and special-function
evaluations are outward rounded.

The decision path is fail-closed.  Every computed enclosure must be finite.
Denominators and square-root domains must be certified positive, the
exponential-integral argument must be certified positive, and branch guards
must be decisive.  An ambiguous comparison is a rejection.  The code uses
the stable factorization \eqref{eq:stable-B}, the positive recursion of
Lemma~\ref{lem:positive-recursion}, the cotangent upper bound
\eqref{eq:cot-tail} with \(100\) terms, and
Corollary~\ref{cor:gaussian-cell}; direct interval subtraction is only the
guarded fallback outside the positive-cosine region.

The endpoint directions in Section~6 are enforced explicitly:
\(B\) is evaluated at \(\rho_+\), the largest \(T\) at \(\rho_-\), and
the \(E_1\) term from the smallest \(T\), which occurs at \(\rho_+\).
The second integral tests several complete upper bounds from
\eqref{eq:cutoff}; taking their finite minimum is sound because no candidate
is merely a bound for a partial tail.

Each worker initializes the constants independently.  Initialization
certifies the isolating bracket for \(\theta_0\), the rational bracket for
\(2\pi\), positivity of the required constants, and the global reduction
\[
 C_E+\frac{3\rhozero}{\sqrt{8124}}<\frac9{20}.
\]

\subsection{Bundle completeness and integrity}

A complete proof object has format
\code{berry-esseen-045-certificate-bundle-v2}.  It contains exactly
\(8122\) ordered shards with the canonical names
\[
 \code{n/0002.json},\ldots,\code{n/8123.json}.
\]
The checker rejects an incorrect scope, missing or additional files,
directories among the shards, renamed files, path traversal, or an incorrect
order.

The manifest records the byte length and SHA--256 digest of each shard.  A
shard is read once; both the parsed object and its digest are derived from the
same byte string.  The checker also recomputes the ordered certificate-set
digest obtained by hashing the ASCII records
\[
 \code{2:hash\_2\string\n},\ldots,
 \code{8123:hash\_8123\string\n}.
\]
The manifest records and the checker recomputes SHA--256 digests for the
numerical core, the public checker, the constructor, the parameter-bank
loader, the one-shot launcher, and both relevant package initializers.  The
parameter-bank manifest digest is retained as construction provenance.  The
constructor and bank remain advice-producing components rather than premises
for a cell inequality: every proposed cell is recomputed by the checker.
These content hashes detect inconsistency within a bundle; they do not
authenticate the repository revision or replace review of the trusted source.

The complete checker performs the following ordered obligations.

\begin{enumerate}[leftmargin=2.2em]
\item Require the specified \code{python-flint} and underlying FLINT versions,
parse the manifest strictly, and require the exact public scope.
\item Validate all \(8122\) entries, source hashes, file names, byte lengths,
shard hashes, the certificate-set digest, and the absence of extra shards.
\item Initialize interval constants and verify the analytic large-\(n\)
reduction.
\item For each \(n\), require an exact gap-free cover of
\([1,\rhozero]\), reconstruct every grid, and recompute every cell.
\item Only after all shards have returned in the canonical order, atomically
write the bundle-local audit record and print \code{ACCEPTED\_FULL}.
\end{enumerate}

Before work begins, the one-shot launcher invalidates any previous accepted
audit.  The checker permits an audit only at the fixed bundle-local path,
rejects symbolic or multiply linked targets, and installs the new record by
atomic replacement.  Every validation or arithmetic failure handled by the
checker prints \code{REJECTED}, returns a nonzero exit status, and leaves no
active accepted global result.  Command-line parsing errors and an external
process interruption also return nonzero, but are reported by the Python
runtime rather than by that validation handler.  A diagnostic
single-shard acceptance is not a certificate of the theorem.

\subsection{Trust boundary}

\begin{table}[htbp]
\centering
\small
\caption{Logical role of the software and data.}
\label{tab:trust-boundary}
\begin{tabularx}{\linewidth}{@{}>{\raggedright\arraybackslash}p{0.31\linewidth}
                              >{\raggedright\arraybackslash}X
                              >{\centering\arraybackslash}p{0.10\linewidth}@{}}
\toprule
Component & Role & Trusted\\
\midrule
Analytic reductions and lemmas in this paper
  & Mathematical derivation of the finite criterion & yes\\
\path{proof/berry_core.py}
  & Outward-rounded realization of equations (20)--(35) & yes\\
\path{proof/check_certificate.py}
  & Strict parsing, coverage, integrity, and complete recomputation & yes\\
Package initializers and \path{construction/verify.py}
  & Faithful one-shot launch, manifest assembly, and result reporting
  & execution\\
CPython, \code{python-flint}, FLINT, and Arb
  & Runtime and interval-arithmetic substrate & yes\\
Linux kernel, Docker engine, CPU, and hardware
  & Host execution substrate outside the container image & yes\\
Certificate constructor and parameter-bank loader
  & Candidate selection, refinement, and checkpointing & no\\
A100 parameter bank and GPU notebooks
  & Floating-point advice used to find efficient rational recipes & no\\
Printed decimals, elapsed time, and worst-cell summaries
  & Human-readable audit information & no\\
\bottomrule
\end{tabularx}
\end{table}

The trusted numerical core and checker form the mathematical
project-specific surface.  The launcher and package initialization path must
also report faithfully when the one-shot command is used; their exact source
bytes are bound into the manifest.  Runtime libraries are execution
assumptions and are pinned by version and binary hashes in the supplied
container.  The host kernel, container engine, processor, and hardware remain
ordinary external execution assumptions: a container fixes userspace, not
those components.  The final checker is launched in a fresh process after
construction.  This separates acceptance from constructor state, but it is
not an implementation-diverse second realization of the analytic formulas.

\subsection{Untrusted construction and reproducibility}

The parameter bank contains \(886{,}998\) dyadic \(\rho\)-cells and six
candidate recipes per cell.  The exact geometry is recovered from integer
depth and index data:
\[
 \left[
 1+\frac{j}{2^d}(\rhozero-1),
 \,
 1+\frac{j+1}{2^d}(\rhozero-1)
 \right].
\]
The bank coordinates for candidate \(q,t_0,\eta,\lambda\) begin on the
dyadic lattice with denominator \(2^{20}\).  The constructor converts them to
exact rationals and may raise \(\eta\) to the exact admissibility boundary
\((\rho_-+1)/q\); consequently, a parameter stored in the final certificate
need not retain denominator \(2^{20}\).  The constructor tries the six
candidates in a fixed order and, if necessary, doubles the proposed number of
grid pieces up to \(4096\).  Every emitted candidate is already evaluated with
Arb, but that constructor-side test is an optimization only; the public checker
recomputes the cell.

Construction is deterministic for fixed sources, configuration, bank, and
runtime.  Each \(n\) is handled in a separate process with no cross-\(n\)
warm start.  Parallelism changes scheduling only.  A completion marker binds
\(n\), the mathematical configuration digest, the imported construction-source
digest, shard digest, and byte length.  Legacy markers from the preceding
format are accepted only while the previously omitted imported sources still
equal their frozen release hashes.  Shards and markers are installed by
atomic replacement, so interruption cannot promote a partial file.  A resumed
run reuses only markers whose content and hashes still match.
The launcher holds an exclusive operating-system lock on the artifact tree for
the entire construction and checking sequence, preventing simultaneous runs
from racing on shards, markers, logs, or the final audit.

The complete Windows execution is:
\begin{lstlisting}[language=bash]
verify.cmd
\end{lstlisting}
It builds the pinned Linux container, constructs or resumes all shards, and
launches a fresh complete check.  The default is six worker processes; output
reports every \(n\) through \(10\) and then every multiple of ten.  The
build targets a pinned \code{linux/amd64} userspace and fixes
Python \(3.12.4\), the Docker frontend and base-image digests, and the hashes
of all installed wheels.  The bank is
mounted read-only, project artifacts are the only writable output, network
access is disabled, Linux capabilities are dropped, and the verifier runs as
a non-root user.

Regeneration of the advice bank is a separate reproducibility experiment.
The output-free notebook
\code{notebooks/build\_parameter\_bank\_colab.ipynb} checks out an immutable
40-hex-character Git commit and requires CPython~3.12 on Linux x86-64 and an
NVIDIA A100 of compute capability~8.0.  The complete Python wheel dependency
closure is version- and SHA--256-locked and forcibly reinstalled; the
Colab-provided CUDA runtime is required to be 12.x, while the exact runtime and
driver versions are recorded rather than fixed by that wheel lock.  Checkpoints bind the build profile, exact
input geometry, parent payload, and saved payload.  A completed bank is copied
to a content-addressed Drive directory and exposed through an atomically
installed pointer only after the copied tree is authenticated.  The builder
publishes a bank only after array shapes, dtypes, byte counts, SHA--256 digests,
exact dyadic coverage, and semantic equality with the committed reference
manifest; it then installs the exact reference manifest bytes.  An A100 is therefore needed to
reproduce the search output, but not to verify a completed certificate.

\begin{proposition}[Checker implication]\label{prop:checker-implication}
Assume the analytic statements used by the trusted numerical core and
faithful execution of the stated trusted stack.  If the public checker
terminates with \code{ACCEPTED\_FULL}, then the hypothesis of
Theorem~\ref{thm:finite-completion} has been verified for every
\(n=2,\ldots,8123\) and every \(\rho\in[1,\rhozero]\).  Consequently,
\[
 \Cstar\le\frac9{20}=0.45.
\]
\end{proposition}

\begin{proof}
Acceptance implies that every required \(n\) occurs exactly once and that
its cells form an exact cover of \([1,\rhozero]\).  On each cell,
outward-rounded evaluation proves \eqref{eq:cell-criterion}.  Their union
gives the hypothesis of Theorem~\ref{thm:finite-completion} throughout the
finite rectangle.  Initialization checks the threshold \(n=8124\), and
Section~2 covers all complementary ranges.
\end{proof}

\section{Exploratory construction of the parameter bank}
\label{sec:exploration}

The parameter advice was obtained by a float64 search on an NVIDIA A100-SXM4
\(80\,\mathrm{GB}\) GPU.  It determines where fine \(\rho\)-subdivision and
dense integration meshes are useful.  Its sampled values are not consumed
by the public checker.

The final bank contains \(886{,}998\) leaves.  Its exact geometry is
summarized in Table~\ref{tab:bank-depth}; the geometry digest is
\[
 \code{bcbe1ee1f583595a929a55043b10e1637a43cc1c76931eddad621c5ef8a3f849}.
\]

\begin{table}[htbp]
\centering
\caption{Final dyadic geometry of the publication parameter bank.}
\label{tab:bank-depth}
\begin{tabular}{@{}rrr@{}}
\toprule
Depth & Final leaves & Percentage\\
\midrule
6 & 430,226 & 48.5036\%\\
8 & 259,884 & 29.2993\%\\
9 & 196,888 & 22.1971\%\\
\midrule
Total & 886,998 & 100\%\\
\bottomrule
\end{tabular}
\end{table}

In the exploratory scoring used to design the bank, the distribution of the
largest sampled score within a cell was as follows.

\begin{table}[htbp]
\centering
\caption{Exploratory float64 scores.  These values are search diagnostics,
not interval bounds.}
\label{tab:exploratory-scores}
\begin{tabular}{@{}lrr@{}}
\toprule
Sampled normalized score & Cells & Percentage\\
\midrule
\(B_{\rm num}<0.4400\) & 886,076 & 99.8961\%\\
\(0.4400\le B_{\rm num}<0.4450\) & 361 & 0.0407\%\\
\(0.4450\le B_{\rm num}<0.4480\) & 467 & 0.0526\%\\
\(0.4480\le B_{\rm num}<0.4495\) & 94 & 0.0106\%\\
\(B_{\rm num}\ge0.4495\) & 0 & 0\%\\
\bottomrule
\end{tabular}
\end{table}

The largest sampled score was
\begin{equation}
 B_{\rm num}^{\max}=0.44820233914501284,
 \manualtag{36}\label{eq:exploratory-maximum}
\end{equation}
with the difficult region concentrated near small \(n\), in particular near
\(n=3\).  The gap to \(0.45\) guided the grid design but has no logical role
in acceptance.

Checkpoint restoration was used during the bank search.  Consequently, the
recorded \(1914.5\)-second active segment in the exploratory notebook is not
an end-to-end reproduction time and is not reported as one.  A fresh
reproduction should report wall-clock time together with which levels were
restored; the complete certificate checker separately records its own elapsed
time in the final audit.

\section{Discussion}

The endpoint estimate in Theorem~\ref{thm:quadratic-tail} changes the
geometry of the verification problem.  A rectangular bound that forgets
\(K(1)=0\) forces the final mesh width to compensate for an \(O(H)\)
artifact.  The regularized identity gives the correct \(O(H^2)\)
contribution, allowing much of the second Prawitz integral to be certified
without an excessively fine endpoint grid.

The partial-fraction expansion and Gaussian-defect identity serve the same
purpose at the other delicate points: they replace subtraction of nearly
equal quantities by positive expressions.  Positivity makes monotonicity
transparent and prevents interval widths from growing merely because two
large overlapping enclosures are subtracted.

The method also illustrates the useful separation between discovery and
verification.  A large GPU search can optimize parameters aggressively
without entering the trusted base.  Once a rational candidate is serialized,
the public checker discards every exploratory score and reconstructs the
mathematical inequality from first principles.  A defective parameter bank
can slow construction or prevent it from finding a certificate; it cannot
cause an invalid cell to pass.

\section*{Data and code availability}

The complete checker, trusted numerical core, certificate schemas,
container specification, parameter-bank provenance, and one-command
reproduction instructions are supplied in the accompanying public
repository.  The certificate shards contain only exact rational data.  The
accepted audit record is written to
\path{artifacts/certificate/verification.json}; the executable checker, not
the human-readable decimal summaries in that record, is authoritative.

\appendix

\section{Exact verification of the threshold
\texorpdfstring{\(n=8124\)}{n=8124}}
\label{app:threshold}

Set
\[
 a=\frac{3162277661}{10^9},
 \qquad
 p=\frac{3141592653}{10^9},
 \qquad
 b=\frac{2506628274}{10^9},
 \qquad
 c=\frac{90133234}{10^6}.
\]
For completeness, define
\[
 A_m(x)=\sum_{k=0}^{m}\frac{(-1)^k x^{2k+1}}{2k+1}.
\]
The alternating-series remainder gives
\(A_{2r+1}(x)<\arctan x<A_{2r}(x)\) for \(0<x\le1\).
Machin's identity \(\pi=16\arctan(1/5)-4\arctan(1/239)\) therefore yields
\[
 \pi>
 L:=16A_7(1/5)-4A_2(1/239)
 =\frac{3367720447004809249690544948}
        {1071978712185331878662109375}.
\]
The entirely rational comparison
\[
 L-p=
 \frac{323056151203315029001}
      {548853100638889921875000000000}>0
\]
proves \(p<\pi\).
Direct integer arithmetic gives
\[
 a^2-10
 =
 \frac{5259630921}{10^{18}}>0,
\]
\[
 2p-b^2
 =
 \frac{495945231}{250000000000000000}>0,
\]
\[
 8124-c^2
 =
 \frac{32175311}{250000000000}>0,
\]
and
\[
 \frac{409732184}{10^9}
 -\frac{a+3}{6b}
 =
 \frac{65406953}{469992801375000000}>0.
\]
Consequently,
\[
 \sqrt{10}<a,
 \qquad
 \sqrt{2\pi}>b,
 \qquad
 \sqrt{8124}>c,
 \qquad
 C_E<\frac{409732184}{10^9}.
\]
Using \(\rhozero=712569/589000\), one more exact subtraction yields
\[
 \frac9{20}
 -
 \left(
 \frac{409732184}{10^9}
 +
 \frac{3\rhozero}{90133234/10^6}
 \right)
 =
 \frac{3121000875001}{3318029676625000000}>0.
\]
This proves
\(C_E+3\rhozero/\sqrt{8124}<9/20\) without floating-point
evaluation.

\section{Correspondence between the manuscript and the trusted core}
\label{app:correspondence}

Table~\ref{tab:correspondence} identifies the principal executable
realizations of the analytic statements.  It is a review map, not an
additional assumption.

\begin{longtable}{@{}p{0.40\linewidth}p{0.54\linewidth}@{}}
\caption{Analytic-to-executable correspondence.}
\label{tab:correspondence}\\
\toprule
Trusted function & Mathematical result\\
\midrule
\endfirsthead
\toprule
Trusted function & Mathematical result\\
\midrule
\endhead
\code{stable\_B\_point}
  & Lemma~4.1 and equation (20)\\
\code{geom\_sum\_upper}
  & Lemma~4.2\\
\code{regularized\_cot\_d\_upper\_point}
  & Lemma~4.3 and equations (21)--(23)\\
\code{kernel\_L\_abs\_upper}
  & Corollary~4.4 and equation (25)\\
\code{rad\_gauss\_diff\_upper}
  & Lemma~4.5, Corollary~4.6, and equations (26)--(28)\\
\code{kernel\_K\_abs\_upper}
  & Lemmas~5.1--5.2 and equations (29)--(30)\\
\code{second\_endpoint\_integral\_upper}
  & Theorem~5.3 and equation (31)\\
\code{second\_integral\_upper}
  & Lemma~5.5 and equation (32), with \(s_\ast\ge1/2\)\\
\code{evaluate\_cell}
  & Lemmas~6.1--6.2, Proposition~6.3, and equations (33)--(35)\\
\code{self\_test\_global\_reductions}
  & Lemma~2.2 and Appendix~A\\
\bottomrule
\end{longtable}

\section*{Certified conclusion}
\label{sec:certified-conclusion}

\ifcertificateaccepted
\begin{finalrecord}
\small
\raggedright
\textbf{Archival verification.}\par\smallskip
The public checker returned \code{ACCEPTED\_FULL} for bundle format
\code{\CertificateBundleFormat}, target \CertificateTarget, and exact scope
\(n=\CertificateScopeStart,\ldots,\CertificateScopeEnd\).  It completed on
\CertificateDate{} using \CertificateRuntime{}, \code{python-flint}
\CertificatePythonFlint, and FLINT \CertificateFlint; the complete check took
\CertificateElapsedSeconds{} seconds and accepted \CertificateTotalCells{}
cells.  The largest reported directed-decimal cell upper bound was
\CertificateWorstUpper{} at \CertificateWorstLocation.

The SHA--256 digest of the accepted manifest is
\par\smallskip
\noindent{\ttfamily\small\CertificateManifestDigest}\par\smallskip
It binds the complete certificate-set digest and every packaged
source digest; the audit record lists the bound digests separately.
\end{finalrecord}

\begin{finaltheorem}
For every \(n\ge1\) and every independent identically distributed sequence
with mean zero, variance \(\sigma^2>0\), and finite third absolute moment,
\[
 \sup_{x\in\R}
 \left|
 \Prob\!\left\{
 \frac{X_1+\cdots+X_n}{\sigma\sqrt n}\le x
 \right\}
 -\Phi(x)
 \right|
 \le
 \frac9{20}\,
 \frac{\E\abs{X_1}^3}{\sigma^3\sqrt n}.
\]
In particular, the universal Berry--Esseen constant satisfies
\[
 \boxed{\Cstar\le\frac9{20}=0.45}.
\]
\end{finaltheorem}

\begin{proof}
The archival bundle-v2 audit records \code{ACCEPTED\_FULL} for target
\(9/20\) and the exact scope \(n=2,\ldots,8123\), after checking the manifest,
certificate-set, and packaged-source digests.  Proposition~\ref{prop:checker-implication}
proves the finite rectangle, and Corollary~\ref{cor:finite-rectangle} proves
every complementary case.
\end{proof}
\else
\begin{resultbox}
\textbf{Archival certificate result to be inserted after the complete run.}
The publication constant is \CertifiedConstant.  Populate the digests, cell
count, worst upper bound and location, elapsed time, completion timestamp, and
environment in
\code{paper/certificate-results.tex} directly from
\code{artifacts/certificate/verification.json}.  First require format
\code{berry-esseen-045-certificate-bundle-v2}, target \code{9/20}, exact scope
\code{2--8123}, the pinned FLINT versions, and matching manifest,
certificate-set, and packaged-source digests; then change
\code{\string\certificateacceptedfalse} to
\code{\string\certificateacceptedtrue}.  No sampled GPU value is suitable
for this record.
\end{resultbox}

\begin{finaltheorem}
If the archival public checker returns \code{ACCEPTED\_FULL} for a verified
bundle-v2 record with target \(9/20\), exact scope \(n=2,\ldots,8123\), and
matching manifest, certificate-set, and packaged-source digests, then the
universal Berry--Esseen constant satisfies
\[
 \boxed{\Cstar\le\frac9{20}=0.45}.
\]
\end{finaltheorem}

\begin{proof}
The accepted audit invokes Proposition~\ref{prop:checker-implication}, and
the conclusion follows from Theorem~\ref{thm:finite-completion} and the
reductions of Section~2.
\end{proof}
\fi

\begin{thebibliography}{99}
\raggedright

\bibitem{BentkusKirsa1989}
V.~Bentkus and K.~Kir\v{s}a,
\newblock Estimates of the closeness of a distribution to the normal law,
\newblock \emph{Lithuanian Mathematical Journal} \textbf{29} (1989),
321--332.

\bibitem{KorolevShevtsova2012}
V.~Korolev and I.~Shevtsova,
\newblock An improvement of the Berry--Esseen inequality with applications
to Poisson and mixed Poisson random sums,
\newblock \emph{Scandinavian Actuarial Journal} \textbf{2012}(2) (2012),
81--105,
\newblock \doi{10.1080/03461238.2010.485370}.

\bibitem{MattnerShevtsova2019}
L.~Mattner and I.~Shevtsova,
\newblock An optimal Berry--Esseen type theorem for integrals of smooth
functions,
\newblock \emph{ALEA, Latin American Journal of Probability and Mathematical
Statistics} \textbf{16} (2019), 487--530,
\newblock \doi{10.30757/ALEA.v16-19}.

\bibitem{Prawitz1972}
H.~Prawitz,
\newblock Limits for a distribution, if the characteristic function is
given in a finite domain,
\newblock \emph{Scandinavian Actuarial Journal} \textbf{1972}(2) (1972),
138--154,
\newblock \doi{10.1080/03461238.1972.10404645}.

\bibitem{Shevtsova2012}
I.~Shevtsova,
\newblock Moment-type estimates with asymptotically optimal structure for
the accuracy of the normal approximation,
\newblock \emph{Annales Mathematicae et Informaticae} \textbf{39} (2012),
241--307,
\newblock Proceedings of the Conference on Stochastic Models and their
Applications, Debrecen, 2011,
\newblock \url{https://ami.uni-eszterhazy.hu/uploads/papers/finalpdf/AMI_39_from241to307.pdf}.

\bibitem{Shevtsova2013}
I.~Shevtsova,
\newblock On the absolute constants in the Berry--Esseen inequality and its
structural and nonuniform improvements,
\newblock \emph{Informatics and Its Applications} \textbf{7}(1) (2013),
124--125,
\newblock \url{https://www.mathnet.ru/eng/ia252}.

\bibitem{ZolotukhinNagaevChebotarev2018}
A.~Zolotukhin, S.~Nagaev, and V.~Chebotarev,
\newblock On a bound of the absolute constant in the Berry--Esseen inequality
for i.i.d. Bernoulli random variables,
\newblock \emph{Modern Stochastics: Theory and Applications} \textbf{5}(3)
(2018), 385--410,
\newblock \doi{10.15559/18-VMSTA113}.

\end{thebibliography}

\end{document}

